\documentclass[11pt]{article}
\usepackage[paperwidth=33.867cm,paperheight=19.05cm,left=1.8cm,right=1.8cm,top=1.6cm,bottom=1.6cm]{geometry}
\usepackage{fontspec}
\usepackage{kotex}
\usepackage{booktabs}
\usepackage{array}
\usepackage{xcolor}
\usepackage{hyperref}
\usepackage{enumitem}
\usepackage{titlesec}
\usepackage{fancyhdr}
\usepackage{amsmath}
\usepackage{float}
\usepackage{caption}
\usepackage{tcolorbox}
\usepackage{tabularx}
\usepackage{colortbl}
\usepackage{setspace}
\usepackage{graphicx}
\usepackage{mdframed}
\usepackage{tikz}
\usepackage{pgffor}

% ---------- Fonts (DGIST report style) ----------
\setmainfont[BoldFont=AppleSDGothicNeo-Bold, ItalicFont=AppleSDGothicNeo-Light]{AppleSDGothicNeo-Regular}
\setsansfont[BoldFont=AppleSDGothicNeo-Bold]{AppleSDGothicNeo-Regular}
\setmonofont{Menlo}

\hypersetup{colorlinks=true, linkcolor=blue!50!black, urlcolor=blue!60!black}

% ---------- Color Palette ----------
\definecolor{dgistnavy}{RGB}{14, 32, 80}
\definecolor{dgistblue}{RGB}{26, 60, 140}
\definecolor{dgistmid}{RGB}{60, 100, 180}
\definecolor{dgistlight}{RGB}{220, 230, 250}
\definecolor{accentgreen}{RGB}{39, 130, 67}
\definecolor{accentorange}{RGB}{210, 100, 20}
\definecolor{accentpurple}{RGB}{100, 50, 150}
\definecolor{accentred}{RGB}{180, 40, 40}
\definecolor{bglight}{RGB}{248, 249, 252}
\definecolor{bggreen}{RGB}{235, 248, 240}
\definecolor{bgorange}{RGB}{255, 247, 235}
\definecolor{bgpurple}{RGB}{248, 240, 255}
\definecolor{bggray}{RGB}{245, 245, 247}
\definecolor{tabhead}{RGB}{26, 60, 140}

% ---------- Header/Footer ----------
\pagestyle{fancy}
\fancyhf{}
\fancyhead[L]{\small\color{dgistblue}\textbf{PROJECT} | DGIST}
\fancyhead[R]{\small\color{gray} 2026년 03월 연구 보고서}
\fancyfoot[C]{\small\color{gray}\thepage}
\renewcommand{\headrulewidth}{0.5pt}
\renewcommand{\headrule}{\hbox to\headwidth{\color{dgistblue}\leaders\hrule height \headrulewidth\hfill}}

% ---------- Section Styles ----------
\titleformat{\section}
  {\large\bfseries\color{dgistnavy}}
  {\colorbox{dgistblue}{\textcolor{white}{\enspace\thesection\enspace}}}
  {0.8em}{}[\vspace{2pt}\color{dgistblue}\hrule height 0.8pt\vspace{4pt}]
\titleformat{\subsection}
  {\normalsize\bfseries\color{dgistblue}}
  {\thesubsection}{0.6em}{}
\titleformat{\subsubsection}
  {\normalsize\bfseries\color{dgistmid}}
  {\thesubsubsection}{0.5em}{}

\setlength{\parindent}{0pt}
\setlength{\parskip}{4pt}
\setlength{\headheight}{22pt}
\addtolength{\topmargin}{-10pt}
\setstretch{1.18}

% ---------- Custom Boxes ----------
\newenvironment{keybox}[2]{%
  \begin{tcolorbox}[colback=#1, colframe=dgistblue!60,
    title={\small\bfseries\color{white}#2},
    fonttitle=\small\bfseries,
    coltitle=white, colbacktitle=dgistblue,
    left=6pt, right=6pt, top=4pt, bottom=4pt,
    boxrule=0.6pt]
}{\end{tcolorbox}}

\newenvironment{insightbox}[1]{%
  \begin{tcolorbox}[colback=bgpurple, colframe=accentpurple!60,
    title={\small\bfseries #1},
    fonttitle=\small\bfseries,
    coltitle=accentpurple, colbacktitle=bgpurple,
    left=6pt, right=6pt, top=4pt, bottom=4pt,
    boxrule=0.6pt]
}{\end{tcolorbox}}

\begin{document}

% ---------- Cover Page (same style, 16:9 size) ----------
\begin{titlepage}
\pagecolor{dgistnavy}
\color{white}
\vspace*{1.0cm}

\begin{center}
{\small\color{dgistlight}\textbf{DGIST / RESEARCH SEMINAR}}\\[0.9cm]

\begin{tcolorbox}[colback=white!8!dgistnavy, colframe=white!30!dgistnavy,
  boxrule=1pt, left=20pt, right=20pt, top=14pt, bottom=14pt]
\centering
{\Large\bfseries\color{white}
Attention Is All You Need}\\[0.45cm]
{\normalsize\color{dgistlight}
Transformer Architecture Brief}
\end{tcolorbox}

\vspace{0.7cm}
{\Large\bfseries\color{dgistlight} 2026년 03월 발표 자료}\\[0.2cm]
{\small\color{gray!60!white} Presentation Deck (Report-Style Theme)}

\vspace{0.9cm}

\begin{tabular}{rl}
\color{dgistlight}\textbf{발표자} & \color{white}Jaeyeong CHOI \\
\color{dgistlight}\textbf{발표명} & \color{white}Attention Is All You Need Presentation \\
\color{dgistlight}\textbf{소속} & \color{white}DGIST \\
\color{dgistlight}\textbf{발표일} & \color{white}2026-03-31 \\
\end{tabular}

\vfill
\end{center}
\end{titlepage}

\pagecolor{white}
\color{black}

\section{연구 개요}
\begin{keybox}{bglight}{핵심 요약}
\textit{Attention Is All You Need}는 RNN/CNN 없이 Self-Attention만으로 시퀀스 변환을 수행하는 Transformer를 제안했고,
현대 LLM(BERT/GPT 계열)의 핵심 아키텍처 기반을 확립했다.
\end{keybox}

\subsection{문제의식}
\begin{itemize}[leftmargin=1.2em]
  \item 장거리 의존성 학습의 어려움
  \item 순차 처리 기반 학습의 병렬화 한계
  \item 대규모 학습 시 계산 비용 급증
\end{itemize}

\section{핵심 방법}
\subsection{Transformer 구조}
인코더-디코더 블록 기반으로 Multi-Head Self-Attention, Feed-Forward, Residual/LayerNorm을 결합한다.

\subsection{핵심 연산}
\[
\mathrm{Attention}(Q,K,V)=\mathrm{softmax}\left(\frac{QK^T}{\sqrt{d_k}}\right)V
\]

\begin{insightbox}{해석 포인트}
\begin{itemize}[leftmargin=1.2em]
  \item $QK^T$로 토큰 간 관련도 계산
  \item $\sqrt{d_k}$ 스케일링으로 학습 안정성 확보
  \item 가중합으로 문맥 표현 생성
\end{itemize}
\end{insightbox}

\section{결과와 영향}
\begin{table}[H]
\centering\small
\begin{tabularx}{\linewidth}{|l|X|}
\hline
\rowcolor{tabhead}\textcolor{white}{\textbf{Category}} & \textcolor{white}{\textbf{Impact}} \\
\hline
Modeling & Recurrence 중심에서 Attention 중심으로 전환 \\
\hline
Efficiency & 병렬 학습 효율 및 확장성 개선 \\
\hline
Legacy & BERT/GPT를 포함한 현대 LLM 구조의 기반 확립 \\
\hline
\end{tabularx}
\caption{핵심 기여 요약}
\end{table}

\section{한계와 시사점}
\begin{minipage}[t]{0.48\linewidth}
\begin{keybox}{bgorange}{한계}
\begin{itemize}[leftmargin=1.2em]
  \item Self-Attention의 $O(n^2)$ 비용
  \item 긴 문맥에서 메모리/지연 증가
  \item 배포 환경에서 효율 최적화 필요
\end{itemize}
\end{keybox}
\end{minipage}\hfill
\begin{minipage}[t]{0.48\linewidth}
\begin{keybox}{bgpurple}{요약}
\begin{itemize}[leftmargin=1.2em]
  \item Transformer 이해는 LLM 이해의 출발점
  \item 성능과 계산 효율을 함께 평가해야 함
  \item Attention 설계 선택이 모델 역량을 좌우
\end{itemize}
\end{keybox}
\end{minipage}

\vfill
\begin{center}
\small\color{gray}\rule{0.5\linewidth}{0.4pt}\\[4pt]
REPORT-STYLE THEME · 16:9 REMAP · 2026-03-31
\end{center}

\end{document}
